\chapter{Background and Related Work}
\label{ch:background}

This chapter presents important context to understand both the problem space and current state of research in passwordless document encryption. Most work addressing real-world problems is interdisciplinary, and this topic is no different. This research overlaps with multiple fields and niche areas including document encryption, usable security, passwordless authentication (FIDO2), and the usability and adoption of passwordless systems. First, a threat analysis of the attack patterns on MFTs which frames the problem space is presented, followed by a literature review of the history of document encryption methods, the challenges of replacing passwords, available open-source encryption tools, and the current state of passwordless authentication provided by FIDO2 and WebAuthn.

\section{Threat Analysis}
\label{sec:threat-analysis}

\subsection{A Pattern of MFT Hacks}
\label{sec:mft-hacks}

From 2020--2025, several mass exploit events compromised millions of confidential documents, including personal health data, government records, and legal files~\cite{kondruss_moveit_2023}. The largest incident resulted in the exposure of 95{,}788{,}491 individuals' personal information~\cite{simas_unpacking_2023}. While each incident had major impacts, perhaps most notable is that this represents a recurring, systematic attack pattern with a specific type of software targeted again and again in similar ways. Capable professional hacking groups---such as the Cl0p hacking group---have repeatedly attacked file transfer platforms including MFT (Managed File Transfer) and FTP (File Transfer Protocol) applications because these systems allow direct access to sensitive files without the need for lateral movement. As Sandra Joyce, a VP at Mandiant noted, ``Mass exploitation of a vulnerability in file transfer software can allow threat actors to efficiently gain access to many organizations' sensitive files without the need for additional lateral movement''~\cite{cisa_cl0p_2023}. The effectiveness of this threat model is demonstrated by its recurrence: the same attack pattern has been exploited against file transfer platforms more than eight times between 2020 and 2025.

Table~\ref{tab:mft_hacks} shows several of the vulnerabilities exploited from 2020--2025, which follow the same threat model. All of the vulnerabilities included are CISA KEV-listed vulnerabilities meaning they are known to be actively exploited ``in the wild.'' The CVSS scores are nearly all high or critical. Many of these were zero-day vulnerabilities discovered and intentionally exploited by capable professional hacking groups including multiple by the same threat actors (e.g., Cl0p).

\begin{table}[htbp]
    \centering
    \renewcommand{\arraystretch}{1.2}
    \setlength{\tabcolsep}{8pt}
    \caption{Actively exploited vulnerabilities in Managed File Transfer (MFT) systems (2020--2025).}
    \label{tab:mft_hacks}
    \begin{tabular}{lllll} 
        \toprule
        \textbf{Date} & \textbf{MFT Provider} & \textbf{CVE ID} & \textbf{CVSS 3.1} & \textbf{CISA KEV} \\
        \midrule
        Dec 2020 & Accellion FTA & CVE-2021-27101 & 9.8 (Critical) & Yes, Added 11/03/2021 \\
        Jan 2022 & SolarWinds Serv-U & CVE-2021-35247 & 5.3 (Medium) & Yes, Added 01/21/2022 \\
        Feb 2023 & Fortra GoAnywhere & CVE-2023-0669 & 7.2 (High) & Yes, Added 02/10/2023 \\
        May 2023 & Progress MOVEit & CVE-2023-34362 & 9.8 (Critical) & Yes, Added 06/02/2023 \\
        Apr 2024 & CrushFTP & CVE-2024-4040 & 10.0 (Critical) & Yes, Added 04/24/2024 \\
        Jun 2024 & SolarWinds Serv-U & CVE-2024-28995 & 8.6 (High) & Yes, Added 07/17/2024 \\
        Dec 2024 & Cleo MFT & CVE-2024-50623 & 9.8 (Critical) & Yes, Added 12/13/2024 \\
        Jul 2025 & Wing FTP Server & CVE-2025-47812 & 10.0 (Critical) & Yes, Added 07/14/2025 \\
        Sep 2025 & Fortra GoAnywhere & CVE-2025-10035 & 10.0 (Critical) & Yes, Added 09/29/2025 \\
        \bottomrule
    \end{tabular}
\end{table}


To illustrate the attack pattern, we will describe details of two of the hacking incidents: Accellion FTA (2020--2021) and MOVEit MFT (2023). These two examples share many similarities to the other vulnerabilities included in Table~\ref{tab:mft_hacks} and are selected because they represent the earliest incident and the largest incident.

\subsection{Accellion FTA Incident}
\label{sec:accellion}

In the Accellion Kiteworks File Transfer Appliances (FTA) incident, which began in December 2020 and extended into 2021, Cl0p chained together four vulnerabilities including zero-day exploits to install a web shell on internet-facing Accellion FTA servers~\cite{cisa_cl0p_2023}. Cl0p used this exploit chain to target federal, state, and local government organizations as well as private organizations across the medical, legal, telecommunications, finance, and energy sectors~\cite{cisa_accellion}. According to data breach reports and legal proceedings, the stolen data included various protected data such as personal health information, Social Security numbers, addresses, emails, and financial information from many large organizations including Bombardier, Jones Day Law Firm, Shell, Kroger, Morgan Stanley, the Reserve Bank of New Zealand, Trinity Health, and the University of California~\cite{msspalert_accellion}.

This incident was significant as it was the first major data breach in what became a series of MFT hacks. The incident occurred in 2020 and impacted diverse sectors from banking and legal to health care, education, and retail. This incident should have been a wake-up call prompting changes in cybersecurity defensive strategy but instead became the first of many similar incidents that would follow over the next five years.

\subsection{MOVEit MFT Incident}
\label{sec:moveit}

The largest incident of this kind was the hack of Progress Software's MOVEit MFT in July 2023. Now a classic example of a third-party exploit leading to data exfiltration, the MOVEit MFT was widely used by many large organizations who were affected. A zero-day SQL injection exploit led to a webshell which gave access to execute operating system commands and exfiltrate files stored in the MFT server~\cite{cisa_cl0p_2023}.

Estimates of affected organizations based on state breach notifications, SEC filings, other public disclosures and Cl0p's website total to 2{,}773 organizations and 95{,}788{,}491 individuals~\cite{simas_unpacking_2023,lyon_moveit}. The organizations with the most impacted individuals are listed in Table~\ref{tab:moveit_impact}.

\begin{table}[htbp]
    \centering
    \renewcommand{\arraystretch}{1.3}
    \caption{Organizations with the most individuals impacted by the MOVEit incident.}
    \label{tab:moveit_impact}
    \begin{tabularx}{\textwidth}{>{\raggedright\arraybackslash}X l >{\raggedright\arraybackslash}p{4.5cm}}
        \toprule
        \textbf{Organization} & \textbf{Impacted Individuals} & \textbf{Industry} \\
        \midrule
        Welltok & 14.76 million & Health care \\
        \arrayrulecolor{gray!40}\midrule\arrayrulecolor{black}
        Maximus & 11.3 million & Government IT Contractor \\
        \arrayrulecolor{gray!40}\midrule\arrayrulecolor{black}
        Delta Dental of California and affiliates & 6.9 million & Dental Insurance \\
        \arrayrulecolor{gray!40}\midrule\arrayrulecolor{black}
        Louisiana Office of Motor Vehicles & 6 million & Government \\
        \arrayrulecolor{gray!40}\midrule\arrayrulecolor{black}
        Alogent & 4.5 million & Banking \\
        \bottomrule
    \end{tabularx}
\end{table}


This incident was significant because it was the largest breach by number of impacted individuals. The total estimated count of 95~million impacted individuals represents 28\%\footnote{$28.65\% = 95{,}788{,}491 / 334{,}233{,}854$, U.S. Census estimated population Jan 1 2023~\cite{us_census_2023}.} of the United States population in 2023~\cite{us_census_2023}. Regardless of the exact numbers it is clear this incident had widespread impacts affecting many ordinary individuals who simply registered their car with the Louisiana Office of Motor Vehicles, or had their dental and healthcare records stored in MFT services from major businesses like Welltok and Delta Dental. Incident investigation found that ``Welltok used the MOVEit Transfer tool for transferring large datasets across the Internet as part of its contracted services with health plans''~\cite{hipaa_welltok_2023}. Similar to the Accellion FTA incident described earlier and the other incidents listed in Table~\ref{tab:mft_hacks}, the MOVEit MFT breach impacted several industries including healthcare, IT, finance, government, and insurance. The common thread across the impacted organizations is the use of large datasets transferred through internet reachable means (i.e., a cloud-based or publicly accessible managed file transfer service).

A notable detail of the MOVEit incident is that the vulnerability patch development and deployment were completed quickly and efficiently. Patches to protect against the MOVEit vulnerability (CVE-2023-34362) were quickly developed once the zero-day issue was identified and were rapidly adopted by affected organizations. One study found that organizations remediated the MOVEit vulnerability roughly 21 times faster than typical remediation rates, reaching in 42 days what would normally take 29 months~\cite{bitsight_moveit_2023}.

Despite the speed with which the patch was developed and implemented across many organizations, millions of individuals were impacted by the data exfiltration that occurred within the brief window when MOVEit servers were vulnerable. This illustrates the key systematic risk of public facing MFTs that do not protect files with per-document granular encryption. In short, the vulnerability patching was fast---essentially a best case scenario for cyber defenders---but because it was a zero-day of public-facing servers and the data was protected in an all-or-nothing encryption approach, it still had a major impact.

\subsection{Summary and Threat Model}
\label{sec:threat-summary}

The significance of the incidents described is that managed file storage options are playing catch up to skilled threat actors instead of designing a better system that can proactively prevent these types of incidents. The repeated pattern of exploits often by the same threat actors targeting the same type of software indicates a systematic problem more than a ``one-off'' issue. Although MFTs encrypt files ``at rest'' all documents are unencrypted and readable by a privileged user session. Once an attacker gains a privileged session on the server---via a zero-day exploit or otherwise---all documents on the server are accessible and can be exfiltrated. While detection, response~\cite{ibm_xforce_mft}, and patching are important reactions to help mitigate the impacts of incidents, in order to fully prevent this attack category, a new defensive approach must be designed and implemented. The current security model fails because it relies on a perimeter defense rather than a zero-trust, per-document protection model. This argument motivates our research. These events highlight the real world need for improved design of document protection systems. In real-world adoption there are always tradeoffs between security and usability; frequently deployability and adoption are the limiting factors for effective security being used in the real world.

Based on the pattern of MFT incidents described above, a general adversarial threat model can be stated. We assume an attacker who has gained initial privileged access to a document repository (e.g., an internet-reachable MFT via a zero-day exploit) and is focused on data exfiltration rather than ransomware or destruction. The threat model does not assume the attacker has local access to individual users' devices or authenticators, rather the attack surface is the server and its stored documents. Table~\ref{tab:threat_mapping} maps each threat to the design requirement that mitigates it and the corresponding security, usability, or deployability benefit from the framework presented in Chapter~\ref{ch:evaluation}.

\begin{table}[htbp]
    \centering
    \renewcommand{\arraystretch}{1.5} 
    \caption{Adversary threat mapping to design requirements and UDS framework benefits.}
    \label{tab:threat_mapping}
    \begin{tabularx}{\textwidth}{
        >{\bfseries\centering\arraybackslash}m{0.75cm}
        >{\raggedright\arraybackslash}X
        >{\raggedright\arraybackslash}X
        >{\raggedright\arraybackslash}X
    }
        \toprule
        & \textbf{Threat} & \textbf{Design Requirement} & \textbf{Benefit(s)} \\
        \midrule
        
        I & Bulk server-side file access & Granular per-document encryption, no shared encryption key & S4 Granular-Protection \\
        
        \arrayrulecolor{gray!40}\midrule\arrayrulecolor{black} 
        
        II & Brute-force password cracking post-exfiltration & Passwordless key derivation & S2 Resilient-to-Guessing \newline U1 Memorywise-Effortless \\
        
        \arrayrulecolor{gray!40}\midrule\arrayrulecolor{black}
        
        III & Unauthorized redistribution of exfiltrated files & Recipient-bound encryption & S3 Resilient-to-Third-Party-Leak \\
        
        \arrayrulecolor{gray!40}\midrule\arrayrulecolor{black}
        
        IV & Server compromise exposing encryption keys & Non-custodial, client-side key derivation, zero-knowledge server & S3 Resilient-to-Third-Party-Leak \\
        
        \bottomrule
    \end{tabularx}
\end{table}


\section{Document Encryption}
\label{sec:doc-encryption}

\subsection{History of Protecting Documents}
\label{sec:doc-history}

Written documents have long held sensitive information and methods to protect that information have developed alongside them. In ancient Mesopotamia, a cuneiform clay tablet dating to 1500 BCE contains the earliest known formula for making pottery glazes. The ancient scribe used cuneiform signs with multiple possible values and truncated symbols in an effort to guard the professional secret~\cite{kahn_codebreakers_1996}. Protecting sensitive information in documents (be it a pottery glaze recipe, legal document, or financial data) remains just as important today, but the methods are often a digital version of the same approach: obscure the contents behind a secret that only the intended reader knows. As digital tools became widespread in the 1990s, the feature to protect a document by encrypting the contents based on a fixed password was developed and added to Microsoft Word and Adobe Acrobat~\cite{adobe_pdf_timeline}. It soon became a standard feature across document file types and operating systems. Today, the encryption works in a similar way as it did when first introduced; however, the technical details of cipher strength and key length have been gradually upgraded. Since the development and distribution of password-protected encrypted documents across Word, PDF, and many operating systems, innovation in protecting documents has lagged.

Efforts have been made by researchers to create simple and cost-effective methods of encrypting individual documents. The most common form of document encryption is password-based AES encryption which is built into Microsoft Word~\cite{microsoft_word_encryption} and PDF processing software like Adobe Acrobat~\cite{adobe_pdf_encryption}. Typically, AES (128 or 256-bit) encryption is used with a single password as the key to protect the entire document. A major problem with this method (used by MS Word and most PDF documents) is that unlike a web-based password, once a password has been set to encrypt a document, it cannot be changed or recovered if forgotten~\cite{microsoft_forgot_password}. One group of researchers discuss the potential implementation of new ciphers for PDF encryption, however the focus of their proposed solution is on minimizing the file size rather than improving the security and usability of the encryption features~\cite{muller_practical_2019}.

\subsection{Protected Documents Today}
\label{sec:doc-today}

Electronic documents containing sensitive data are used across nearly every industry and profession. Examples include health records, legal counsel, financial data, trade secrets, and employment contracts. These documents frequently must be protected, especially before being shared externally. The PDF format has become a de facto standard for encrypting and transmitting sensitive documents. M\"uller et al.\ (2019) document that PDF encryption is widely deployed: companies like Canon and Samsung use it in document scanners to protect sensitive information, IBM offers PDF encryption services for wrapping confidential data including images, and healthcare organizations use encrypted PDFs to transfer protected health records. Some organizations have adopted gateway systems that automatically encrypt outgoing email messages as password-protected PDF attachments, with the password communicated over a second channel such as SMS~\cite{muller_practical_2019}. The prevalence of PDF encryption across these use cases underscores both the importance of the format and the urgency of addressing its limitations, which are examined in detail in Chapter~\ref{ch:evaluation}.

These sensitive documents are often stored and shared using cloud services like managed file transfer (MFT) applications, as discussed in the threat analysis in Section~\ref{sec:threat-analysis}. The focus of this research is document protection in the sense of confidentiality: ensuring that only the intended recipients can read the contents of a document. However, confidentiality without usability or deployability is impractical because such a system will not be used in real-world scenarios; thus the requirements are security, usability, and deployability~\cite{bonneau_quest_2012,borisov_sok_2021}. Related topics such as document integrity, digital signatures, and content modification detection---which primarily address the integrity aspect of cybersecurity---are outside the scope of this work.

\subsection{Document Encryption Hacking}
\label{sec:doc-hacking}

Research on document encryption security has focused primarily on demonstrating weaknesses in password-based protection. M\"uller et al.\ in their 2019 CCS paper~\cite{muller_practical_2019} trace a pattern of researchers cracking PDF passwords and ensuing patches from 2001 through 2020. In 2001 Komulainen et al.\ performed a security analysis of encrypted PDFs and pointed out risks of 40-bit key length encryption. Sklyarov et al.\ demonstrated brute-force tools to crack PDF encryption at DEF CON 9 in the same year. This led Adobe\footnote{Adobe owned and managed the PDF standard from inception until 2008 when it became an ISO standard 32000 which is managed today by the PDF Association.} to increase the key length from 40 to 128 bits for RC4 encryption in the PDF version 1.4. In 2008, Sklyarov et al.\ returned and found security issues which enabled brute-force attacks in the latest PDF version 1.7. Adobe again updated the key derivation in the PDF standard. In 2013, Danczul et al.\ showed methods for brute-forcing encrypted PDF passwords across different processors including GPUs~\cite{danczul_cuteforce_2013}. This research found that nearly any password-based document encryption can be cracked with ``little overhead'' given enough time and the proper computational resources~\cite{danczul_cuteforce_2013}. In 2015, August at al.\ measured the time to brute force a PDF password based on the password's length. In 2017, Stevens et al.\ demonstrated a method to crack PDF passwords within seconds by downgrading to old algorithms and using tools like pdf2john. Some commercial tools such as Elcomsoft apply these techniques and sell the ability to break encrypted documents which are used by law enforcement agencies and others. Results from prior research have also found methods to break XML encryption. In 2019, M\"uller et al.\ demonstrated a novel solution for breaking the confidentiality of a protected PDF by adding exfil content at the bottom of the PDF which executes when the document is opened~\cite{muller_practical_2019}. The most recent and relevant research article on breaking PDF encryption emphasized that there are likely additional security flaws in document encryption, but limited research has been performed in the area~\cite{muller_practical_2019}. Novel methods of breaking the encryption of password protected documents have been developed and published as recently as November 2024~\cite{hranicky_experimental_2024}.

Beyond document encryption, several research efforts have demonstrated weaknesses with current file security measures. In 2021, researchers published their findings of two new vulnerabilities allowing undetected content modification on certification-signed 2.0 PDFs~\cite{rohlmann_breaking_2021}. These findings resulted in new CVEs being created for several software vendors who in turn released off-cycle patches to remediate the security issues found. The researchers found that 24 of the 26 commonly used PDF viewer applications they tested were vulnerable to attack~\cite{rohlmann_breaking_2021}. Security researchers have also uncovered vulnerabilities in the digital signature process used to sign and verify PDF~\cite{rohlmann_breaking_2021}, Microsoft Office (OOXML), and OpenDocument files. Novel attacks such as content masking and modification of signed PDF files have also been discovered and reported. Microsoft's Rights Management Services, which enforces access controls on Word documents across Azure and Office 365, was also found vulnerable to attacks that could completely strip protection from view-only documents~\cite{grothe_rms_2016} (see MSRC Case 33210).

In summary, despite decades of ongoing vulnerabilities and patching, the typical approach to document encryption of deriving the offline key based on a user-selected password to protect the file remains largely unchanged.

\subsection{Design Research}
\label{sec:doc-design-research}

Design research on document protection methods has primarily focused on maintaining privacy and integrity. For example, CryptPad is an open-source anonymous document editing software for which the design process has been analyzed. Efforts to improve the integrity of the PDF specification include strengthening the PDF trust chain~\cite{tullsen_weak_2022} and employing the zero-trust file methodology of Content Disarm and Reconstruction (CDR)~\cite{dubin_cdr_2023}. Research to maintain offline document integrity resulted in a Microsoft Word plugin that tracks document changes and provides integrity and non-repudiation, but does not handle confidentiality~\cite{shatnawi_maintaining_2017}. Further related work has addressed serializing document encryption enforcement and investigated a blockchain method to protect documents. While design research on protecting documents is active there is relatively little design work focused on the underlying confidentiality mechanism to protect a document's contents.

\subsection{Usability Gap}
\label{sec:usability-gap}

Regarding the usability of encrypted documents, very little work has been done to directly assess the usability of current document protection methods. However, extensive usability research has been performed on passwords~\cite{bonneau_quest_2012}, passwordless technology~\cite{alaca_comparative_2021-1}, and encrypted email~\cite{borisov_sok_2021}, which each presents important connections to the problem of usability in document encryption. Passwords are the incumbent solution for document protection but have key differences when compared to website login passwords used for user authentication. Unlike web-based authentication, password-protected documents cannot throttle guess attempts, cannot recover or reset a forgotten password, and offer no mechanism for revoking access once the password has been shared. These differences mean that password-based document encryption provides less security and less usability than passwords used for web authentication (a method that itself has well-documented limitations). The implications of these limitations and the broader quest to replace passwords are discussed in Section~\ref{sec:quest-passwords}.

\section{The Quest to Replace Passwords}
\label{sec:quest-passwords}

Considerable research has focused on the quest to replace passwords with new authentication methods that provide improved security and usability benefits. In 1999, the study ``Why Johnny Can't Encrypt'' evaluated PGP through a lab study and found that PGP's encryption system ``[was] not usable enough to provide effective security for most computer users, despite its attractive graphical user interface, supporting our hypothesis that user interface design for effective security remains an open problem.''~\cite{whitten_why_1999}. This study became a landmark paper in usable security as it highlighted the need to study usability and user acceptance of security systems in order to understand user adoption. Passwords specifically have known usability problems of increased cognitive load which leads to weak security practices including password reuse and the selection of simple passwords. The reader is likely familiar from lived experience of this. As researchers note, ``With every new user account and every new password cognitive load is increasing so that usability is often preferred over security by users: For example, users often choose the same password across accounts, keep an insecurely stored written record or choose unsecure dictionary passwords'' (e.g., Adams et al.\ 1997, Johnson and Grawemeyer, 2011, Wash et al., 2016).

Despite the security weaknesses and usability flaws of passwords, however, their use persists. Another landmark paper in usable security zeroed in on this question: Why have so many password-replacement concepts not caught on? Their systematization of knowledge on the topic involved a formalized evaluation of passwords and password-replacement methods against a set of benefits grouped into three categories: usability, deployability, and security~\cite{bonneau_quest_2012}. This methodology and formal evaluation approach is adapted in this work by applying the UDS benchmarking to the document encryption context in Chapter~\ref{ch:methodology}. The authors~\cite{bonneau_quest_2012} find that while password-alternative methods often provide superior security and/or usability, no system provides equal or superior benefits in deployability. Deployability includes a low per-user cost to set up and operate a system. A highly deployable system is compatible with existing technology and can be implemented and used by anyone without having to pay a fee. The authors conclude that deployability is a major reason why no method has gained traction to replace traditional passwords for user authentication. In many respects, 2-factor-authentication is to balance the security, usability, and deployability tradeoffs that exist between different authentication methods.

In a later article continuing the analysis, Herley and Van Oorschot (two of the authors of \emph{The Quest to Replace Passwords} study) write ``The `password replacement problem' is both underspecified (in vagueness of specific goals or concrete requirements) and over-constrained (it is simply impossible to find a single solution addressing all security and usability needs in all scenarios). Regarding the latter, rather than seeking a silver bullet, we must consider best-fit solutions'' after acknowledging that ``[passwords] are likely to persist for some time'' they point out that ``there are scenarios where passwords are not the best-fit. To determine these cases and find suitable alternatives requires (a) clearly and specifically identifying the requirements; and (b) devising a practical, reliable methodology to evaluate, score, and compare competing alternatives against these requirements. The research avenue suggested is to gain better insight into actual threat likelihoods and actual harms experienced when things fail. An important step is to replace unsupported opinions by factual evidence and real-world data''~\cite{herley_research_2011}. This research is grounded in this approach starting with an analysis of ``actual harms experienced when things fail'' as described in Section~\ref{sec:threat-analysis}, followed by a set of requirements and evaluation method presented in Chapter~\ref{ch:methodology} and~\ref{ch:evaluation} which is then assessed in a lab study in Chapter~\ref{ch:user-study}.

While the quest to replace passwords has focused primarily on user authentication, the same usability-deployability-security tensions apply to document encryption. Encrypted documents present a related but distinct model with two important differences. First, an encrypted file often must be shared with another user to access, where a user's login should not be shared with others for access. Second, a document or file must be self-sufficient, meaning it functions standalone from any server, while user authentication often depends on connection to a server either through the internet or directly through user login. These two differences shift the usable security model for document encryption compared to user authentication: multiple parties are expected to access the same resource, and the file must function independently of a server.

The practical consequences of these differences are visible in how organizations handle sensitive documents today. Even government agencies with formal data protection requirements rely on password-based file encryption with manual password communication. For example, the U.S.~Department of Education instructs individuals submitting documents with personally identifiable information (PII) to ``encrypt [and email] the file and send a second email\ldots\ with the password to decrypt that file''~\cite{doe_encryption}. This two-email process---one for the file, one for the password---illustrates the usability and security limitations that motivate this research: the sender must create, remember, and securely transmit a password through a separate channel for every encrypted document shared and the process is prone to interception, unauthorized sharing, and human error.

The challenges of protecting shared documents have many similarities with secure email---in fact, sensitive documents are often sent as encrypted email attachments---thus the usability research performed in this domain has important implications for protected documents. Like email, document file formats were not originally designed with security in mind, and most encryption features have been bolted on after the fact. A systematization of knowledge on secure email found that regular users often do not encrypt their communications, and that usability is not the only obstacle to adoption~\cite{borisov_sok_2021}. Garfinkel et al.\ found that half of surveyed users did not use encrypted email because they did not know how, while the remainder indicated they did not think it was necessary or that the effort would be wasted~\cite{garfinkel_email_2005}. More broadly, users do not perceive significant risk, lack knowledge about effective ways to mitigate risk, and lack confidence in their ability to use secure systems effectively~\cite{borisov_sok_2021}. These findings parallel the document encryption context directly: if users do not understand how to encrypt a document, do not perceive the risk of leaving it unencrypted, or do not feel confident they can do it correctly, adoption will not follow regardless of the technical security properties of the tool.

Beyond secure email and password replacement research, broader cybersecurity research has identified self-efficacy---a belief in one's own ability to enact cybersecurity-related skills---as a key predictor of security behavior adoption. An evidence review by the European Union Agency for Cybersecurity (ENISA) found across studies of Protection Motivation Theory and Theory of Planned Behaviour that self-efficacy was ``a reliable, moderately strong predictor of cyber-security intention and behaviour'' outweighing threat-based factors~\cite{enisa_cybersecurity_2019}. A systematic review of 174 studies reinforced the centrality of self-efficacy in cybersecurity behavior research~\cite{borgert_systematic_2024}. Self-efficacy has also been shown to predict perceived ease of use, which in turn predicts perceived usefulness of a technology~\cite{zhang_self_efficacy_2017}. When users feel confident in their ability to use a tool, they perceive it as easier and more worthwhile to use. Importantly, response efficacy must accompany self-efficacy, meaning users must believe that they are capable of using the tool and that the tool will effectively protect them in order for user acceptance and adoption to occur~\cite{enisa_cybersecurity_2019}. In short, interventions that enable users to feel confident in using a security tool are more likely to increase user adoption than interventions that emphasize the severity of threats. These findings inform the qualitative themes explored in the user study in Chapter~\ref{ch:user-study} and underscore that the problem space is not simply a technical encryption challenge but relies on human-centric psychological factors.

In summary, the literature on password replacement and security behavior demonstrates that achieving real-world adoption requires addressing usability and deployability alongside security, and that human-centric factors such as self-efficacy and perceived effort are as important as the technical properties of the system itself~\cite{bonneau_quest_2012}.

\section{Open Source Encryption Tools}
\label{sec:oss-tools}

Deployability of a system is most likely when the system can be implemented without fees or significant material costs. Proprietary systems with licensing fees often struggle to gain widespread adoption due to the financial disincentive. As such, open source software (OSS) is an appealing area to consider for this problem space as the software is typically provided free of use and can be installed within various architectures without resulting in licensing fees. There are several open source data encryption standards and software libraries. This section will share a few open source tools that are particularly relevant to the research performed.

Age is an open-source, modern file encryption tool designed by Filippo Valsorda~\cite{age_repo}. A formal specification of age is defined as part of the Community Cryptography Specification Project~\cite{age_spec}. Designed in part to replace GPG for file encryption, age has small explicit keys and UNIX-style composability. Age uses authenticated encryption (ChaCha20 as cipher and Poly1305 as message authentication code) in a streaming chunk design to provide confidentiality and integrity by default and support large files. Since version 1.3.0 age includes support for post-quantum schemes (via ML-KEM-768 and X25519).

Initially developed as a Go library to perform file encryption, the age specification has been implemented in multiple programming languages and supports plugins and integrations with several encryption management tools. A Rust implementation (rage) and TypeScript implementation (typage) exist alongside plugins for YubiKeys and other identity authenticators. Age is designed to support multiple recipients where each recipient can independently decrypt the file. In its default mode, age encrypts to either a passphrase (password) or public key string(s). Both of these modes require users to manage and exchange cryptographic material out-of-band. Age is intentionally focused on file encryption and does not handle key exchange or secure messaging, leaving those problems to other tools. Notably, the typage browser implementation supports passkey-based encryption via the WebAuthn PRF extension, discussed in Section~\ref{sec:webauthn-prf}, which provides an alternative to password-based and manual key exchange for file encryption.

Picocrypt~\cite{picocrypt_repo} (now maintained as Picocrypt NG) is a lightweight, open-source file encryption tool with strong cryptographic features. While not used directly in this work, it is an alternative option for open-source file encryption. Picocrypt also relies on password-based key derivation focusing on the file encryption problem instead of the key distribution or access control problems relevant to sharing encrypted documents.

OpenTDF~\cite{opentdf_spec} is the open-source implementation of the Trusted Data Format which is a standard of the U.S.\ Director of National Intelligence~\cite{dni_tdf} that enables attribute-based access control (ABAC) to be cryptographically tied to a data object. ABAC enables more granular access control than the all-or-nothing access control of traditional password-based file encryption. The role of OpenTDF in future extensions of this work is discussed in Section~\ref{sec:opentdf}.

These tools primarily address the file encryption process and do not fully address how a user can authenticate---prove their identity---without the factor of something you know (e.g., password, secret key). This requires a separate set of protocols discussed in the next section.

\section{Passwordless Authentication and FIDO2}
\label{sec:fido2}

While MFTs continue to be hacked and sensitive documents are protected either individually by passwords or all together by cloud based storage, there has been vast improvement in recent years to develop and standardize passwordless authentication protocols.

\subsection{Hardware-based Authentication}
\label{sec:hw-auth}

One of the most well-known weaknesses of passwords is that they are often exploited in phishing attempts because users assume the phishing site is legitimate. To remedy this, early work on passwordless authentication focused on creating phishing-resistant forms of user authentication. Many modern devices now include a secure hardware enclave---such as a Trusted Platform Module (TPM) on Windows and Linux devices and Secure Enclave on Apple devices---a hardware-based security feature that can generate and store cryptographic keys isolated from the device's main processing components. Hardware tokens such as YubiKeys provide a similar secure hardware enclave to calculate random numbers which is used to generate and store cryptographic keys. The biometrics in modern devices (e.g., a fingerprint or facial scan) work with the on-device secure enclave to provide hardware-bound secrets and remove the password and memory aspect from the user authentication. A hardware token works similarly as a portable physical asset that is separate from the computer device. These methods of hardware-based authentication secure user login by transferring the work of key creation and management from the user to a dedicated computational device. This means users can authenticate their identity without relying on memory (passwords, security questions) or basic 2FA (text, email code). This work to build standardized hardware-based authentication---motivated by an effort to have phishing resistant authentication---in many ways led to the FIDO2 standardization (for instance, Yubico, Google, Apple, and Microsoft who included trusted compute environments in their devices have been important contributors to FIDO2 standards).

\subsection{Passwordless Authentication Outside of FIDO2}
\label{sec:other-passwordless}

Various passwordless authentication schemes have been proposed; some are now defunct or are proposed in academic research but are not widely used in practice today. These approaches include:

\begin{itemize}
    \item Hardware Authenticator Binding
    \item Pico
    \item Let's Authenticate \emph{(Automated Cryptographic Authentication for the Web with Simple Account Recovery)} and \emph{Let's Authenticate: Automated Certificates for User Authentication}
    \item Mozilla Persona
    \item TUSH-Key: Transferable User Secrets on Hardware Key
\end{itemize}

However, none have achieved the industry adoption or standardization momentum of FIDO2.

\subsection{FIDO2}
\label{sec:fido2-core}

FIDO2 is an open standard for passwordless authentication developed and maintained by the FIDO (Fast IDentity Online) Alliance, whose membership includes major technology companies such as Google, Microsoft, and Apple. Unlike earlier passwordless authentication schemes, FIDO2 has substantial industry backing that has driven widespread support across browsers, operating systems, and common websites. This gives FIDO2 the potential to overcome the deployability barriers that have historically prevented password replacement mechanisms from gaining traction.

As of 2026, this potential is being realized at scale. The FIDO Alliance's \emph{State of Passkeys 2026} report, based on a survey of 11{,}000 consumers across ten countries, found 90\% of consumers are now familiar with passkeys and 75\% have enabled them on at least some accounts. In workforce environments, 68\% of organizations are deploying, piloting, or rolling out passkeys for employee authentication~\cite{fido_state_of_passkeys_2026}. In total, there are an estimated 5 billion active passkeys globally~\cite{fido_state_of_passkeys_2026}. These figures represent a significant shift from even two years prior and suggest that FIDO2 has begun to cross the adoption threshold from a novel technology to mainstream use.

At its core, the FIDO2 process is made up of two standardized protocols: CTAP2 (Client to Authenticator Protocol)~\cite{fido_ctap_2025} and WebAuthn (W3C's Web Authentication API)~\cite{w3c_webauthn_l3}. Together, these protocols shift authentication from the knowledge factor (``what you know'') to the possession factor (``what you have''). There is no shared secret between the server and client. FIDO2 is resistant to phishing, replay attacks, and server-side credential breaches because unique keys are created for each authentication ceremony and the secret is never stored or shared with the server. The user's secret is derived on their authenticator which can be either a platform authenticator (e.g., Windows Hello, Apple Face ID) where the keys are typically stored in the device's secure hardware enclave, or a roaming authenticator (e.g., YubiKey) which is device-independent. Users authenticate locally to the authenticator using a PIN or biometric factor such as a fingerprint or facial scan. The CTAP protocol defines communication between the client browser and the authenticator, while WebAuthn defines the API through which the browser communicates with the server (called the relying party). Figure~\ref{fig:fido2-arch} illustrates this architecture.

\begin{figure}[htbp]
    \centering
    \includegraphics[width=5in]{fido2-architecture}
    \caption[FIDO2 architecture.]{FIDO2 architecture. The authenticator creates and stores the private key, communicating with the client via the CTAP protocol. The client communicates with the relying party server via the WebAuthn API. The server begins the process and stores the public key.}
    \label{fig:fido2-arch}
\end{figure}

By integrating CTAP and WebAuthn into a standard with backing from major technology vendors, FIDO2 offers a standardized passwordless authentication protocol with broad support across the modern internet infrastructure. Deployment is promising but not yet universal. A recent analysis of 111 websites found that FIDO2 passkey implementations generally adhere to the FIDO Alliance Design Guidelines approximately 77\% of the time, with higher-ranked websites exhibiting more consistent passkey user experiences and closer alignment with the guidelines~\cite{ramat_passkey_2025}. Jannett et al.\ found steady growth in passkey adoption from 2021--2025 across the internet but that many sites had security flaws in their implementation of WebAuthn~\cite{jannett_webauthn_2026}. In summary, FIDO2 is gaining significant traction but is not yet without implementation challenges.

\subsection{WebAuthn PRF Extension}
\label{sec:webauthn-prf}

FIDO2 and WebAuthn provide a robust authentication ceremony with many benefits, but as they are built for user authentication they do not on their own provide file encryption. The remaining question is: how do you derive encryption keys without a password or server-stored key? The WebAuthn PRF extension addresses this gap.

WebAuthn PRF is an optional extension in the Web Authentication Level 3 specification (Section 10.1.4) where the relying party provides a salt to the authenticator via WebAuthn which the browser hashes with a context string before passing it to the authenticator, which then combines the hashed salt with its internal key via HMAC-SHA-256 to produce a deterministic 32-byte string which is returned to the relying party~\cite{w3c_webauthn_l3}. A pseudo-random function (PRF) produces output that is deterministic for a given input but computationally indistinguishable from random, which makes the output string suitable for use as symmetric key material. The W3C specification directly names symmetric key encryption as the motivating example for the inclusion of the PRF extension~\cite{w3c_webauthn_l3}. The underlying protocol mechanism is the CTAP2 hmac-secret extension, which provides access to the hardware-backed HMAC-SHA-256 computation. The WebAuthn PRF extension acts as a standardized browser API for web applications to access this capability~\cite{corbado_prf_2026}. The process of WebAuthn PRF is described in Figure~\ref{fig:prf-flow}.

\begin{figure}[htbp]
    \centering
    \includegraphics[width=5in]{webauthn-prf-flow}
    \caption[WebAuthn PRF extension flow.]{WebAuthn PRF extension flow across the four required layers (authenticator, operating system, browser, and web application) showing the process for deriving a symmetric key from the authenticator's internal secret.}
    \label{fig:prf-flow}
\end{figure}

While standard WebAuthn operates as a challenge-response for user authentication, the PRF extension enables a deterministic output that can be used as a symmetric encryption key. The security benefits of this model include a unique key per user per website (each site has a unique salt), hardware-bound derivation (the internal key never leaves the authenticator's secure element), tamper resistant (the authenticator is a trusted compute environment), non-custodial (the relying party never receives the authenticator's secret), and phishing resistant (WebAuthn scopes credentials to the relying party's origin)~\cite{yubico_prf_extension,bitwarden_prf_2024}.

To understand why PRF improves security beyond existing two-factor authentication, consider how a hardware security key functions differently in each model. When a YubiKey is used for traditional 2FA, it presents a second gate that must be passed through before the server shares encrypted data to the requesting user. If an attacker bypasses 2FA through a server-side compromise or session hijacking, the data's protection falls entirely to the password because the data encryption was based on the password and the 2FA served as an additional gate, not a part of the encryption derivation process itself. PRF fundamentally changes this model by making the passkey part of the encryption process itself rather than an access gate in front of it. The encryption key is derived directly from the authenticator's hardware-bound secret combined with the relying party's salt: no password is involved. This means that even if an attacker compromises the server, they obtain only encrypted ciphertext with no feasible path to the decryption key, because the decryption key exists only when the authenticator is physically present and the user completes the WebAuthn ceremony~\cite{bitwarden_prf_2024,yubico_prf_extension}. In short, 2FA protects \emph{access} to encrypted data; PRF protects the \emph{encryption itself}.

The PRF extension has seen early implementation in use cases where security and a non-custodial model is of primary importance. Password managers including Bitwarden and 1Password use PRF to enable vault decryption without a master password~\cite{bitwarden_prf_2024}. Digital wallets like cryptocurrency wallets and digital identity wallets, such as the wwWallet project, use PRF to derive encryption keys for non-custodial credential storage~\cite{yubico_prf_extension}. These early adopters demonstrate PRF's viability as a realistic solution for non-custodial encryption but focus on personal data vaults. To our knowledge no prior work applied the WebAuthn PRF scheme to file encryption with multi-party sharing as is needed for secure document management.

Support for WebAuthn PRF requires orchestration across four layers: the hardware authenticator, the operating system, the browser, and the web application. A failure at any layer means the extension is non-operable (Figure~\ref{fig:prf-flow}). As of mid-2026 support is maturing rapidly but is not yet universal. Chromium-based browsers (Chrome, Edge) support PRF on all major desktop platforms. Windows Hello added hmac-secret support in February 2026 (via cumulative update KB5077181), resolving a longstanding gap on Windows~\cite{corbado_prf_2026}. Firefox 148+ supports PRF with Windows Hello on both credential creation and authentication~\cite{corbado_prf_2026}. Safari added PRF support for iCloud Keychain passkeys (platform authenticators) in Safari 18, and extended support for roaming authenticators (security keys) in an April 2026 update.

In short, WebAuthn PRF enables secure symmetric encryption where the secret is tied to the hardware-based passwordless authenticator. This means a passkey can serve not only as a login credential but also as a source for enabling passwordless, non-custodial file encryption. This is the technical foundation for the system design presented in Chapter~\ref{ch:system-design}.

\subsection{Passkeys: Synced vs.\ Device-Bound}
\label{sec:passkey-types}

While FIDO2 and its underlying CTAP2 protocol can operate with a variety of authenticators, one of the most common is a passkey which is a term FIDO2 presents as ``the high level, end-user centric term for a FIDO2/WebAuthn Discoverable Credential''~\cite{passkeys_dev_terms}. ``Like `password', `passkey' is a common noun'' which is ``designed to be used without additional login challenges''~\cite{passkeys_dev_terms}. Each passkey is created and stored by a passkey provider which is responsible to manage the passkey and is often called a password manager or credential manager. The passkey provider may be native to the operating system (first party, e.g.\ ``Windows Hello'' on Windows, ``Apple Passwords'' on macOS and iOS, and ``Google Password Manager'' or ``Samsung Pass'' on Android devices) or may be added via a browser extension or other tool that intercepts WebAuthn requests and negotiates with the platform authenticator (third party, e.g.\ ``Bitwarden'', ``1Password'', etc.).

There are two types of passkeys: synced and device-bound. There are important security-usability-deployability tradeoffs between synced and device-bound passkeys. A synced passkey is shared across a user's devices typically through a backend cloud-sync by the passkey provider which enables the user to restore passkeys from a backup and have the passkeys available on different devices. For example if a user created a synced passkey on their phone using Google Password Manager, Google would sync the passkey to the cloud so the user could access and login with the same passkey on their laptop computer if they had authenticated to Google Password Manager on the laptop. Synced passkeys provide usability benefits of easy recovery from loss, cross-platform sharing, and ease of use but have a security tradeoff as the passkey credential is not hardware-bound, can be used on multiple devices, and is stored and synced by the passkey provider instead of fully non-custodial. A device-bound passkey is bound to a single authenticator such that the credential cannot leave the device. A FIDO2 security key like a YubiKey or a PRF-enabled passkey are examples. Device-bound passkeys are also called single-device passkeys. When considering portability---the ability to move between devices---it is important to understand the distinction between platform and roaming authenticators. A platform authenticator is built-in to a user's device~\cite{passkeys_dev_terms} while a roaming authenticator is a physically separate device attached by USB, NFC, or Bluetooth which can be used to sign into multiple devices. A roaming authenticator is often called a security key. Examples of roaming authenticators include a mobile smartphone, a YubiKey, and more. A roaming authenticator may create either a device-bound or a synced passkey depending on the configuration and the user's choice. This means there are $2\times 2$ form factor options for the creation of passkeys, as summarized in Table~\ref{tab:passkey_form_factors}.

\begin{table}[htbp]
    \centering
    \renewcommand{\arraystretch}{1.5} 
    \caption{Security and usability differences between FIDO2 Passkey form factors.}
    \label{tab:passkey_form_factors}
    \begin{tabularx}{\textwidth}{
        >{\bfseries\raggedright\arraybackslash}m{2.8cm} % Vertically centered, bolded row headers
        >{\raggedright\arraybackslash}X               % Synced column
        >{\raggedright\arraybackslash}X               % Device-bound column
    }
        \toprule
        % Note: \textbf{} is omitted for "Authenticator" because the entire column is set to bold above
        Authenticator & \textbf{Synced} & \textbf{Device-bound} \\
        \midrule
        
        Platform & Created by the built-in platform authenticator. The credential can be discovered and used on multiple devices via cloud sync. & Created by a single device's platform authenticator. Only usable on that specific hardware device. \\
        
        \arrayrulecolor{gray!40}\midrule\arrayrulecolor{black} 
        
        Roaming & Created on a separately attached device, but the credential can be discovered and used on multiple devices. & Created on a separately attached device. The credential cannot leave the authenticator, but the physical authenticator can be moved to login on multiple devices. \\
        
        \bottomrule
    \end{tabularx}
\end{table}


Functionally speaking a synced-platform and synced-roaming passkey have similar security-usability-deployability characteristics as they both afford the ability to authenticate from multiple devices based on the syncing with the passkey provider. When the passkey is synced, the device where the credential was first created doesn't change the security posture since any device with access to the passkey provider's vault can access the credential. The security and usability tradeoffs between these FIDO2 form factors have been evaluated by researchers using the Bonneau et al.\ UDS framework and are discussed in the following section.

\subsection{Security and Usability Evaluations of FIDO2 Authentication}
\label{sec:fido2-evals}

Security-focused researchers have examined FIDO2, CTAP2, and WebAuthn in a variety of methods including formal security proofs, local attacks, and other novel attacks. For example, Barbosa et al.\ perform provable cryptographic security and privacy analysis of FIDO2's WebAuthn and CTAP2 protocols~\cite{barbosa_provable_2021}. Guan et al.\ use ProVerif to formally verify security assumptions about FIDO2~\cite{guan_formal_2022}. Kuchhal et al.\ evaluate the security of FIDO2 deployments across 1{,}000 common websites and find that compromised clients (e.g., with malware) pose a practical threat to FIDO2 deployments~\cite{kuchhal_fido2_2023}. Tarun and Seamons present novel local attacks against FIDO2 including malicious browser extension attacks which in some cases ``present more risk than passwords alone''~\cite{tarun_security_2024}. These works generally find FIDO2 has a strong security framework but needs some improvements and is dependent on a non-compromised local execution environment.

Usable security researchers have examined the benefits of passwords vs.\ FIDO2 hardware tokens and passwords vs.\ FIDO2 synced vs.\ FIDO2 device-bound passkeys. In their paper ``Is FIDO2 the Kingslayer of User Authentication?'' Lyastani et al.\ performed a large-scale user study of passwords vs.\ FIDO2 hardware tokens finding that lay users were generally ``very satisfied'' when replacing text-based passwords with single-factor security keys; however, users lacked mental models to understand how security keys worked and had concerns about the suitability for everyday use, authenticator revocation, account recovery, and overall trustworthiness of the FIDO2 system. The authors conclude that FIDO2 has the potential to be the kingslayer of passwords despite significant shortcomings~\cite{lyastani_fido2_2020}:
\begin{quote}
We find it astonishing that users accept [FIDO2 1FA] authentication with security keys so strongly despite these shortcomings. The main reason could be that the disadvantages and weaknesses of text-based passwords have become so obvious and overwhelming for users that they are looking for a technology that can free them from this burden. In summary, we find that there is still a gap between the users' concerns and what the current status-quo of FIDO2 1FA provides.
\end{quote}

In terms of the Bonneau et al.\ U-D-S framework, Lyastani et al.\ find that FIDO2 hardware tokens provide nearly all the benefits with four and a half gaps: \emph{Nothing-to-Carry}, \emph{Easy-Recovery-from-Loss}, \emph{Negligible-Cost-Per-User} (partial), \emph{Server-Compatible}, and \emph{Resilience-to-Theft}. They differentiate the benefits which are universal to FIDO2 and which are dependent on the authenticator device. This study focused specifically on hardware tokens which lack the benefits of \emph{Nothing-to-Carry}, \emph{Negligible-Cost-Per-User} (partial), and \emph{Resilience-to-Theft} due to the physical hardware device. Platform FIDO2 authenticators (passkeys) offer the same passwordless FIDO2 protocols and benefits without a separate physical device that must be carried, purchased, and could be stolen. As such platform passkeys (both device-bound and synced) provide additional benefits when compared to FIDO2 hardware tokens (which are included in the device-bound roaming authenticator category).

B\"uttner and Gruschka performed a comparison of passwords, device-bound passkeys (both platform and roaming), and synced passkeys using the U-D-S framework presented by Bonneau et al.\ and found that ``the security of synced passkeys is mainly concentrated in the passkey providers''~\cite{buttner_passkey_2025}. Notably, B\"uttner and Gruschka evaluate device-bound passkeys as providing all security benefits in the framework while synced passkeys lack the \emph{No-Trusted-Third-Party} security benefit as a breach of a passkey provider (e.g., iCloud Keychain, Google Password Manager, 1Password) could lead to account compromise thus making the passkey provider a ``single point of failure'' for synced-passkeys~\cite{buttner_passkey_2025}. However, the same security benefit of device-bound passkeys being tied to specific hardware and thus not having a trusted third party is why device-bound passkeys lack the usability benefit of \emph{Easy-Recovery-from-Loss}.

W\"ursching et al.\ conducted a lab study comparing user acceptance of FIDO2 platform and roaming authentication on smartphones by having users sign into a mock website through either TouchID or a YubiKey via NFC. Results showed SUS scores for both platform and roaming authenticators $>$90 indicating excellent usability and acceptance scores for both methods similarly very high. The authors note that ``our participants want to use passwordless authentication on smartphones when possible'' and that platform authentication had higher acceptance than roaming authentication on smartphones. They find that the account type (e.g.\ banking vs.\ email vs.\ streaming service) changes user adoption preferences based on the perceived sensitivity of the online account~\cite{wursching_fido2_2023}. Similar to the other works discussed, users were hesitant to adopt FIDO2 due to concerns about account recovery, credential sharing, and confusion about how multi-device access works for platform authenticators. Notably a countertrend was observed of some users rejecting passwordless authentication for their most sensitive accounts trusting their memory over a passwordless process: ``I consider the [YubiKey] to be secure. However, for important accounts, I would rather rely on my memory and that the login does not get into the wrong hands.''~\cite{wursching_fido2_2023}. This finding correlates with the observation by Lyastani et al.\ that users lack a clear mental model of how FIDO2 operates so while users generally see FIDO2 as ``secure'' when faced with sensitive data they may fall back to trusting a password-based mechanism.

Beyond controlled lab studies, field studies have uncovered practical barriers to FIDO2 adoption and deployment. Owens et al.\ compared a smartphone as FIDO2 roaming authenticator to passwords for user login to a fictitious bank website over two weeks and found that while participants saw the FIDO2 method as more secure they found it less usable than passwords with concerns about the need for a physical device, account recovery, and setup challenges~\cite{owens_user_2021}. Ciolino et al.\ collected diary entries over a week and found users preferred SMS codes to FIDO2 security keys with users reporting uncertainty about the security benefits of security keys~\cite{ciolino_twominds_2019}. Lassak et al.\ interviewed 28 CISOs and authentication managers and found barriers to adoption of account recovery, friction, technical issues, regulatory requirements, and security culture~\cite{lassak_why_2024}. These results suggest that the conceptual usability and security benefits identified through academic frameworks or lab studies may not fully translate to practical benefits in sustained real-world use.

Amer et al.\ investigated whether technical explanations about FIDO2's benefits would reduce user's reluctance to adopt FIDO2. They found that explicitly communicating security benefits did not influence participant's intention to adopt FIDO2 and did not overcome usability concerns. Notably they found participant's intention to adopt FIDO2 was correlated with factors like self-efficacy and risk perception~\cite{amer_explaining_2024}. This finding reinforces the self-efficacy literature discussed in Section~\ref{sec:quest-passwords} that user adoption may be driven more effectively by building users' confidence in their ability to use the security tool than emphasizing its security benefits. For example, Reittinger et al.\ found a significant increase in passkey adoption when applying digital nudges at appropriate touchpoints during a randomized controlled trial with 3{,}680 participants on a commercial healthcare platform~\cite{reittinger_passkey_2025}.

In summary, FIDO2 presents strong security and usability benefits, though notable concerns exist around account recovery, user mental models, and feasibility in practical day-to-day use. Platform authenticators generally show usability advantages over roaming authenticators; but account recovery and sustained user adoption remain open challenges across all FIDO2 form factors. Of these challenges, account recovery stands out as the most salient concern across the literature and is discussed in the following section.

\subsection{FIDO2 Account Recovery and Sharing}
\label{sec:fido2-recovery}

Researchers consistently identify account recovery as ``one of the most controversial aspects'' of FIDO2 authentication~\cite{buttner_passkey_2025}. Lyastani et al.\ found that recovery was ``a predominant [user] concern\ldots\ in line with prior user study results on 2FA with security keys''~\cite{lyastani_fido2_2020}. W\"ursching et al.\ observed similar results in their user study~\cite{wursching_fido2_2023}. Farke et al., Kepkowski et al., and Lassak et al.\ showed the same concerns about account recovery held in corporate and enterprise contexts~\cite{farke_security_2020,lassak_why_2024}. The primary reason given by studies by Das et al.~\cite{das_johnny_2018} and Ciolino et al.~\cite{ciolino_twominds_2019} for users not using 2FA security keys are that users are afraid of losing their key and locking themselves out of the account or that users see no need to protect the account with an extra security effort~\cite{farke_security_2020}. Kepkowski et al.\ in an applied usability study with 118 enterprise professionals found account recovery was selected by 60\%+ of participants as a top concern with using FIDO2 in their enterprise setting~\cite{kepkowski_challenges_2023}.

A few approaches to device-bound FIDO2 credential recovery have been proposed and evaluated in research but have not been widely adopted as the FIDO Alliance generally presents synced passkeys or registering a second security token as the solution for account recovery~\cite{kunke_evaluation_2021,buttner_passkey_2025}. B\"uttner and Gruschka note that there are levels of access and sharing for passkeys from device-bound which is single-user, single-device to synced which is typically single-user but multi-device, shared which enables multiple users to access the same key, and exported where the credentials are bundled and removed from the passkey-provider~\cite{buttner_passkey_2025}. Each level of sharing increases the access level for account recovery and third-party risk. Various users and organizations may adjust their permitted access levels for passkey sharing based on their risk tolerance for third-party breaches and their needs for efficient account recovery.

\begin{figure}[htbp]
    \centering
    \includegraphics[width=5in]{passkey-sharing-levels}
    \caption[Access levels of passkey sharing.]{Access levels of passkey sharing, from device-bound to synced, shared, and exported. Adapted from B\"uttner and Gruschka (2025)~\cite{buttner_passkey_2025}.}
    \label{fig:passkey-sharing}
\end{figure}

Ultimately the choice to use a device-bound or synced passkey rests with the user as there are security vs.\ usability tradeoffs with both options. Synced passkeys introduce provider trust by shifting security to the passkey provider, but enable multi-device portability and account recovery while device-bound passkeys have no third party risk but lack easy recoverability. Thus the weaknesses of device-bound and synced passkeys are largely complementary. A system that supports both authenticator types---allowing users to choose based on their security needs---could approach the ideal balance of usability, deployability, and security. This possibility is explored in the system design presented in Chapter~\ref{ch:system-design}.
